
\subsection{Equivalent critical continuation conditions}
\label{r:sec:terminal-characterization}

The projected acceleration and critical dissipation are linked by
an exact identity. Combined with the spatial energy estimate, it
characterizes continuation through \(T\) in terms of several
equivalent quantities.
Define
\[
 \mathcal B_T=\int_0^T\norm u_{\dot H^{3/2}}^2\dd t.
\]
Projection of the material momentum law gives the exact identity
\begin{equation}
 \mathcal B_T
 =\nu^{-2}\norm{\Pp(D_tu-f)}_{L^2(0,T;\dot H^{-1/2})}^2.
 \label{r:eq:projected-response-identity}
\end{equation}
The identity follows from \(\Pp(D_tu-f)=\nu\Delta u\) and Plancherel.
It holds on each strict interval and passes to the nonnegative terminal
integrals, including the value \(+\infty\).

For the stated smooth datum and source, the following assertions about
a classical forced history on \([0,T)\) are equivalent:
\begin{equation}
 \begin{gathered}
 u\in\mathscr I_T;\qquad
 \mathcal B_T<\infty;\qquad
 \Pp D_tu\in L^2(0,T;\dot H^{-1/2});\\
 \widehat C_1[u]\in L^1(0,T);\qquad
 \sup_{t<T}\int_0^t\mathcal P_{1/2}(s)\dd s<\infty.
 \end{gathered}
 \label{r:eq:terminal-equivalent-statements}
\end{equation}
Here \(\widehat C_1\) is the original completion formula evaluated on
this forced velocity, as in \eqref{r:eq:original-form-completions}.
The source has a finite \(L^2\dot H^{-1/2}\) norm, so subtraction of
\(\Pp f\) preserves the projected-acceleration condition.
Equation \eqref{r:eq:low-completion-source-shift} identifies the
completed-coordinate condition with \(\mathcal B_T<\infty\).

For the remaining implications, first suppose \(\mathcal B_T<\infty\).
At spatial height one, differentiation and incompressibility give
\[
 |\mathcal P_1|
 \le\norm{\nabla u}_3\norm{\nabla u}_2\norm{\nabla u}_6
 \le K_1\norm u_{\dot H^{3/2}}\norm{\nabla u}_2\norm{\Delta u}_2.
\]
Young's inequality in the forced energy balance gives
\[
 \frac{\dd}{\dd t}\norm{\nabla u}_2^2+\nu\norm{\Delta u}_2^2
 \le\frac{2K_1^2}{\nu}\norm u_{\dot H^{3/2}}^2
               \norm{\nabla u}_2^2
       +\frac2\nu\norm{\Pp f}_2^2.
\]
Thus Gronwall bounds the \(H^1\) norm uniformly through \(T\).
Local classical existence in \(H^1\), with the prescribed smooth source,
has a lifespan bounded below by its \(H^1\) bound and a source bound.
Restarting sufficiently near \(T\) extends the history past \(T\).
Uniqueness identifies that extension on its overlap.
Higher spatial differentiation and the momentum recurrence then supply
every fixed spatial and temporal Sobolev norm on the closed horizon.
Equivalently, the explicit higher-energy estimates of
Section~\ref{r:sec:higher-spatial-source} perform this reconstruction.
This proves \(u\in\mathscr I_T\).
The reverse implication follows by selecting the finite spatial coordinate.

Finally, the work \(W_{1/2}=\ip u{\Lambda\Pp f}\) is uniformly bounded
by kinetic energy and a prescribed source norm. The exact critical balance
is
\[
 E_{1/2}(t)+\nu\int_0^t\norm u_{\dot H^{3/2}}^2\dd s
 =E_{1/2}(0)+\int_0^t\mathcal P_{1/2}\dd s
                  +\int_0^tW_{1/2}\dd s.
\]
A uniform upper bound for the signed production integral controls both
nonnegative terms on the left, including \(\mathcal B_T\).
Conversely, membership in \(\mathscr I_T\) gives the absolute production
bound \eqref{r:eq:critical-production}. This proves all the equivalences.

These equivalences characterize full-interval regularity of a given forced
solution. For the global solution of Part~II,
Corollary~\ref{r:thm:supplied-forced} supplies all these conditions and every
finite-rectangle bound.
